%% =============================================================================
%%  Etapa 5 — Intertexto
%% =============================================================================
\chapter{Intertexto}\label{cap:intertexto}

\begin{objetivo}
Ouvir o Antigo Testamento por trás de Lucas 6: os textos que o trecho cita,
evoca ou pressupõe, lidos em hebraico e aramaico, e comparar com a antiga
tradução siríaca. Perguntar, em cada caso, se há citação, alusão ou apenas
eco, e o que muda na leitura \parencite{beale2007}.
\end{objetivo}

\section{Isaías 61 e o programa de Nazaré}

Em Nazaré, Jesus lê \rb{Is 61:1–2} e declara: \enquote{Hoje se cumpriu esta
Escritura} (\rb{Lc 4:18–21}). Na Septuaginta, o texto reúne três palavras que
reaparecem em Lucas~6:

\begin{citacao}{\rb{Is 61:1–2} LXX (trechos)}
\gr{εὐαγγελίσασθαι πτωχοῖς ἀπέσταλκέν με \ldots\ παρακαλέσαι πάντας τοὺς πενθοῦντας}
\end{citacao}

\noindent \gr{πτωχοῖς} (os pobres, \rb{Lc 6:20}), \gr{παρακαλέσαι}
(consolar, cf.\ \gr{παράκλησις}, \rb{6:24}) e \gr{πενθοῦντας} (os que choram,
cf.\ \gr{πενθήσετε}, \rb{6:25}). As felicidades soam como o anúncio de Is~61
aplicado aos discípulos; os ais, como o avesso desse anúncio.

\section{Isaías 65:13–14: servos e \enquote{vós}}

Isaías 65 contrapõe duas sortes: \enquote{meus servos} comerão, beberão,
se alegrarão e cantarão; \enquote{vós} tereis fome, sede, vergonha e
gritareis de dor. A estrutura em pares e a segunda pessoa são muito
próximas de Lucas.

\interlinear[titulo={Is 65:13–14 · Códice de Leningrado (WLC)}]{dados/biblia/is65_13-14.hbo.tsv}

\begin{pergunta}
Faça uma tabela com as quatro oposições de Is 65:13–14 ao lado das quatro
de Lc 6:20–26. Quais coincidem em vocabulário? Quais só na ideia?
\end{pergunta}

\section{Salmo 1: o macarismo hebraico}

\hb{אַשְׁרֵי} (\tl{ʾašrê}) é um substantivo plural em estado construto:
literalmente \enquote{felicidades de\ldots!}, uma exclamação. A Septuaginta
traduz \gr{μακάριος ἀνήρ}.

\interlinear[titulo={Sl 1:1 · WLC}]{dados/biblia/sl01_01.hbo.tsv}

\section{Daniel 7:13: \enquote{como um filho de homem}}

\enquote{Por causa do Filho do Homem} (\rb{Lc 6:22}) é o título que Jesus
mais usa para si mesmo nos Evangelhos. Seu pano de fundo mais citado está em
aramaico, na visão de \rb{Dn 7:13}: \arc{כְּבַר אֱנָשׁ} (\tl{kəbar ʾĕnāš}),
\enquote{como um filho de homem}, isto é, uma figura humana que recebe domínio
do Ancião de Dias.

\interlinear[modo=leitura, titulo={Dn 7:13 · WLC (aramaico bíblico)}]{dados/biblia/dn07_13.arc.tsv}

\begin{nota}
Em Lc 6:22, sofrer \enquote{por causa do Filho do Homem} liga o destino dos
discípulos ao de Jesus, que também será rejeitado (\rb{Lc 9:22}) e depois
exaltado (\rb{Lc 22:69}), como a figura de Daniel.
\end{nota}

\section{A Peshitta: Lucas 6 em siríaco}

A Peshitta é a tradução siríaca clássica do NT, de uso contínuo nas igrejas
de tradição siríaca. O siríaco é um dialeto aramaico: próximo da língua de
Jesus, mas \emph{não} é o texto original dos Evangelhos, que foram escritos em
grego. Sua leitura ajuda a ouvir as palavras num idioma semítico.

\interlinear[titulo={Lc 6:20b e 6:24 · Peshitta}]{dados/biblia/lc06_20b-24.syr.tsv}

\paragraph{Observações.} \sy{ܛܘܒܝܟܘܢ} (\tl{ṭūbaykōn}) é
\enquote{felicidades vossas}, com o mesmo tipo de construção exclamativa do
hebraico \hb{אַשְׁרֵי}. O \enquote{ai} é \sy{ܘܝ} (\tl{way}); e a consolação dos
ricos é \sy{ܒܘܝܐܐ} (\tl{buyāʾā}), da raiz \enquote{consolar}. A BP traduz
esse texto, o que explica suas escolhas próprias na Etapa~1.

\begin{alerta}[A conferir]
O texto siríaco foi transcrito das edições disponíveis em linha (texto de
Gwilliam/Urmia) e a análise morfológica é provisória. Confirme com uma
edição impressa e com um léxico (Payne Smith; SEDRA) antes de usar em
conclusões.
\end{alerta}

\section{Ecos no Novo Testamento}

\begin{ebtabela}{@{}l >{\raggedright\arraybackslash}X@{}}
\EBcab{Texto} & \EBcab{Eco de Lc 6:20–26}\\\midrule
\rb{Tg 2:5} & Deus escolheu os pobres deste mundo para serem ricos na fé e herdeiros do reino\\
\rb{Tg 5:1–6} & \enquote{Chorai e lamentai, ó ricos}: uma série de advertências aos ricos\\
\rb{1Pe 4:14} & \enquote{Se sois insultados pelo nome de Cristo, felizes sois} (\gr{ὀνειδίζεσθε \ldots\ μακάριοι})\\
\rb{Lc 16:25} & a inversão entre o rico e Lázaro; Lázaro é \enquote{consolado}\\
\rb{Lc 1:53} & o Magnificat: famintos saciados, ricos despedidos vazios\\
\end{ebtabela}

\begin{pergunta}
Tiago parece conhecer uma tradição próxima de Lucas 6. Que diferenças de
ênfase você percebe entre \rb{Tg 5:1–6} e \rb{Lc 6:24–26}?
\end{pergunta}
